% =========================================================================
% General Conclusion and Perspectives
% =========================================================================
\chapter*{General Conclusion and Perspectives}
\addcontentsline{toc}{chapter}{General Conclusion and Perspectives}
\markboth{Conclusion and Perspectives}{}

During this end-of-studies project we were able to design and implement
\textbf{BViral Control Center}, a four-pillar AI platform that helps short-form video
creators on TikTok, Instagram Reels and YouTube Shorts answer the three questions that
matter the most before they hit publish: \emph{will this clip perform, are the captions
readable, is the visual quality up to par?}

The Pre-Posting Predictive Agent fuses CLIP visual embeddings, Librosa audio descriptors,
Whisper + MiniLM transcript embeddings, and Llama-3 generated semantic features into a
473-dimensional PCA-projected representation, predicts a virality score with an XGBoost
regressor that achieved R\textsuperscript{2} = 0.83 on the held-out test set, and explains
the prediction with SHAP TreeExplainer values. The Dynamic Auto-Captioning Engine pairs
\texttt{faster-whisper} with word-level timestamps and a safe-font-size algorithm that
guarantees zero overflow on portrait-mode video, reducing word error rate to 6.7\,\% and
removing the 27\,\% of clipped captions observed on competitor tools. The AI Video
Enhancement Pipeline assesses video quality with OpenCV, runs classical filters,
Real-ESRGAN super-resolution and GFPGAN face restoration, and produces a manifest of
profiles (mobile, tablet, desktop, 4K) with associated PSNR, SSIM and LPIPS metrics for
quality assurance. The AI Studio dashboard, built on React/Vite with a Fastify API,
PostgreSQL with Drizzle ORM, Supabase Storage, BullMQ on Redis and a vendored Next.js
video editor, orchestrates the three pillars end-to-end with multi-platform OAuth
(YouTube, TikTok PKCE, Meta) and scheduled publishing.

This report has detailed in its three chapters the various stages of building this
solution: the state of the art with the analysis of existing tools (VidIQ, TubeBuddy,
CapCut AI, Descript) and the requirement specification; the conceptual study with the UML
diagrams, the dataset description and the system architecture; and finally the realisation
with the four AI pillars, the technology choices, the development environment and the
user-facing interfaces.

\bigskip

While coding this project, several technologies, libraries and design patterns were
internalised: PyTorch and \texttt{float16} GPU inference, the BullMQ queue architecture for
long-running AI jobs, PostgreSQL Row Level Security with Drizzle declarative policies, the
TikTok PKCE OAuth flow, the ASS subtitle override syntax, and the Real-ESRGAN tiled
inference strategy that keeps VRAM bounded.

\bigskip

\noindent\textbf{Future perspectives.}
This work has achieved most of its initial objectives, but several promising directions
remain open and should be explored in future iterations.

\begin{itemize}
  \item \textbf{Real-time streaming analysis.} Today the predictive agent runs on uploaded
        files. A WebRTC ingest path that streams a few seconds of footage at a time would
        allow the AI Coach to give a live score \emph{while the creator is recording},
        opening the door to a viewfinder-style coaching overlay.
  \item \textbf{Native mobile application.} A React Native client would let creators trigger
        the entire pipeline directly from their phone, reusing the same Fastify API and the
        same Zod schemas through the generated TanStack Query client.
  \item \textbf{Federated learning.} The current Pillar~I model is trained on a centralised
        TikTok scrape. A federated training scheme in which each creator's private engagement
        data updates a personalised head on top of a frozen multimodal backbone would let
        every creator benefit from a personalised virality model without exposing their data.
  \item \textbf{Multi-platform social-network integration.} Today Snapchat is declared in
        the platform enum but not implemented. Adding Snapchat (Snap Kit) and X (formerly
        Twitter) publishers would extend the reach of the scheduler.
  \item \textbf{Caption translation.} The captioning engine could be extended with a
        post-Whisper translation step (NLLB-200 or M2M-100) so that a single recording can
        be exported to multiple languages with localised on-screen text.
  \item \textbf{Diffusion-based enhancement.} Real-ESRGAN does well for super-resolution
        but a diffusion model such as Stable Video Diffusion XL or Stable Cascade Video
        could deliver perceptually superior frames at the cost of more compute. A
        diffusion-vs-Real-ESRGAN A/B test integrated into the dashboard would be the natural
        next experiment.
  \item \textbf{Continuous evaluation.} Once a critical mass of post-publication analytics
        has been collected through the platform, the Pillar~I regressor can be retrained
        regularly on real first-party data, and a calibration plot (predicted score vs
        observed engagement) can be surfaced to the user as a trust signal.
\end{itemize}

\cleardoublepage
